\chapter*{Úvod}
\addcontentsline{toc}{chapter}{Úvod}

Zdravý životní styl a sledování osobních návyků hrají v dnešní době důležitou roli. Na trhu sice existuje množství aplikací pro fitness, výživu či spánek, ale nutnost používat více nástrojů současně bývá nepřehledná a demotivující. Proto vznikla myšlenka mobilní aplikace \enquote{LifeMeter}, která tyto klíčové aspekty sjednocuje do jednoho uživatelsky přívětivého prostředí.

\section{Cíl práce}
Cílem maturitního projektu je návrh a vývoj aplikace \enquote{LifeMeter} pro operační systém Android (verze 10 a vyšší). Požadované funkce zahrnují:

\begin{itemize}
    \item \textbf{Záznam fyzických aktivit:} Možnost přidávat aktivity s informacemi jako poloha nebo poznámka.
    \item \textbf{Sledování výživy:} Přidávání zkonzumovaného jídla, doplňků stravy, sledování kalorií, makroživin a vitamínů.
    \item \textbf{Sledování spánku:} Záznam délky spánku s možností přidání poznámek, případně možnost pracovat s kvalitou spánku. 
    \item \textbf{Správa uživatelů:} Přihlašování bude probíhat pomocí emailu se zabezpečeným ověřením. 
    \item \textbf{Offline:} Aplikace bude umožňovat offline režim se synchronizací dat při obnovení připojení. 
    \item \textbf{Notifikace:} Upozornění ke cvičení a spánku.
    \item \textbf{Integrace s jinými aplikacemi:} Aplikace bude umožňovat propojení s vybranou aplikací/platformou.
\end{itemize}

\section{Struktura práce}
Text práce je rozdělen do navazujících částí. 
% TODO: přidat popis jednotlivých kapitol